\section{Results}
\label{sec:results}

All results are reported under Leave-One-Participant-Out Cross-Validation
(LOPO-CV): in each fold, one participant is held out entirely as the test set,
and all preprocessing, SMOTE generation, and model training are performed
exclusively on the remaining participants.
This is the strictest available cross-participant evaluation protocol and
directly simulates deployment on a driver the system has never encountered.
The primary performance metric is the area under the ROC curve (AUROC), which
measures the model's ability to rank impaired windows above safe ones
independently of any operating threshold; calibration quality is additionally
assessed through Expected Calibration Error (ECE)~\cite{Guo2017} and Brier
score.
All aggregate metrics are accompanied by 95\% confidence intervals obtained
via non-parametric bootstrap ($N = 2\,000$ resamples).

Not all participants contribute equally to this evaluation.
Those with fewer than five positive windows in their held-out fold provide
insufficient signal for reliable AUROC estimation and are designated as
\emph{audit drivers}.
For this group, the relevant safety question shifts from discrimination to false
alarm control: rather than AUROC, we report the false positive rate (FPR) at a
fixed operating threshold, quantifying how often the model raises an impairment
alert for a driver who committed no safety-critical error.

Before reporting main results, we additionally present a systematic ablation
over the two key windowing hyperparameters: the lookback duration $L$ and the
labelling horizon $H$.
The configuration selected from this ablation is used for all results that
follow.

\subsection{Main Results}

Table~\ref{tab:main_results} reports AUROC, ECE, and Brier score for KinTCN
and the two baselines: logistic regression (LR) and gradient-boosted trees
(XGBoost~\cite{Chen2016}), both trained on the same window-level summary
statistics (mean, standard deviation, and maximum per signal).
KinTCN with Platt calibration achieves an AUROC of
$\mathbf{XX} \pm XX$ $[XX, XX]$, compared to $XX \pm XX$ for XGBoost and
$XX \pm XX$ for the LR baseline.
The gap between KinTCN and the tree-based baseline is consistent
across participants, with KinTCN achieving higher per-participant AUROC in
$XX$ out of $XX$ evaluable drivers (Figure~\ref{fig:per_driver_auroc}).
Calibration follows the same pattern: KinTCN achieves an ECE of $XX$
and a Brier score of $XX$, against $XX$ and $XX$ for XGBoost,
confirming that the Platt recalibration step yields reliable probability
estimates rather than merely well-ranked scores.

\subsection{Statistical Significance}

Statistical significance of the improvement over XGBoost is assessed using the
Wilcoxon signed-rank test on per-participant AUROC values, with a two-sided
alternative and $\alpha = 0.05$.
The test yields $W = XX$, $p = XX$, indicating that KinTCN's advantage over
the gradient-boosted baseline is [significant / not significant] at the chosen
level.
The mean per-participant gain is $\Delta\mathrm{AUROC} = +XX \pm XX$,
with KinTCN exceeding XGBoost in $XX$ of $XX$ evaluable drivers.

\subsection{Permutation Feature Importance}

To understand which kinematic signals drive the model's decisions, we compute
permutation feature importance: for each signal in turn, its values across the
test windows are randomly shuffled and the resulting drop in per-participant
AUROC is recorded.
A larger drop indicates a signal the model relies on more heavily.
Figure~\ref{fig:perm_importance} reports the mean drop and its standard
deviation across drivers and permutation repeats.
[DESCRIBE FINDINGS: e.g., steering torque and lateral acceleration account for
the largest drops, consistent with the literature on early collision signatures;
speed contributes moderately; steering wheel angle is least informative.]
This ranking is stable across participants, suggesting that the model has
learned a consistent kinematic profile of impairment rather than overfitting
to participant-specific driving styles.

\subsection{Safe Driver Audit}
\label{sec:audit}

Audit drivers are participants whose held-out fold contains fewer than five
positive windows, meaning they made very few or no safety-critical errors
during the study session.
For this group, AUROC is uninformative; the relevant safety question is instead
whether the model raises false alarms at an unacceptable rate for an attentive
driver.
For each audit driver, a dedicated LOPO fold is run---holding that driver out
entirely and training on all others---and the false positive rate (FPR) at a
fixed operating threshold of $0.5$ is measured on their negative windows.
Across the $XX$ audit drivers, the mean FPR is $XX \pm XX$
(Table~\ref{tab:audit}), and the mean calibrated impairment probability
assigned to their safe windows is $XX$, well below the alert threshold.
These figures indicate that drivers who drive safely throughout the session are
unlikely to trigger false alarms, a key requirement for any deployable
monitoring system.

\subsection{Ablation: Lookback Duration and Predictive Horizon}

The predictive window is governed by two hyperparameters: the lookback
duration $L$, which controls how much kinematic history the network can
attend to, and the labelling horizon $H$, which determines how far into the
future errors are aggregated.
Figure~\ref{fig:ablation_grid} reports AUROC on a grid of $L \times H$ values,
with the gap $G$ held fixed.
[DESCRIBE FINDINGS: e.g., performance increases with $L$ up to $XX$\,s and
then plateaus, suggesting that kinematic precursors beyond that horizon carry
little additional information; horizon $H = XX$\,s offers the best
trade-off between coverage and temporal precision.]
The chosen configuration ($L = XX$\,s, $H = XX$\,s, $G = XX$\,s) sits at
the performance peak and is used for all results reported above.
